% p1-acknowledgments

\section{P1 Acknowledgments}

Renormalization
  Dmitrii Vasilev, Stergios Pellis, Scott Olsen

  GOLDEN BRIDGE (Trinity-Pellis Bridge) Programme --- Paper 2 of 3

  Trinity  $S^{3}$AI   ---   Flos          Aureus    Research     Programme      $\cdot$   Zenodo    record
  10.5281/zenodo.19227877

          Epistemic status (mandatory reading). This paper is a mechanism-exploration
          document. It does not derive Standard Model coupling constants from first
          principles, and no claim to that effect is made or implied. The results of Paper 1
          --- symbolic compressibility of $\varphi$-structured coupling patterns --- are taken as
          empirical motivation only; they are not proof of a physical mechanism. The
          theoretical frameworks developed here --- the symbolic RG model, the candidate
          operator framework, the discrete-scale-invariance connection --- are conjectural
          constructions whose status is scaffolding, not established theory. All
          mathematical statements are clearly labeled as either (a) known exact results
          imported from the literature, (b) model-internal derivations whose assumptions
          are explicitly stated, or (c) speculative proposals awaiting formal derivation or
          experimental test. Sections that cross the boundary into conjecture are marked
          [CONJECTURAL]. The goal is to identify a research programme with clear
          success and failure criteria, not to announce a discovery.

  Abstract
  Paper 1 of the GOLDEN BRIDGE (Trinity-Pellis Bridge) programme established that a suite of coupling-constant
  combinations from the Standard Model and near-critical quantum systems admit
  unusually short symbolic descriptions when encoded in a basis containing the golden
  ratio $\varphi$ = (1+5)/2. Whether this compressibility reflects a genuine physical mechanism
  or is an artefact of number theory and selection bias remains open. The present paper
  asks a sharper question: are there known, derivable mechanisms in quantum field
  theory and statistical mechanics that could cause $\varphi$-structured patterns to appear in
  renormalized coupling constants?

  Three candidate lines of evidence are examined. First, the Zamolodchikov E8 integrable
  field theory --- established in two seminal papers published in Adv. Stud. Pure Math. 19
  (1989) 641--674 (Project Euclid) and Int. J. Mod. Phys. A 4 (1989) 4235 --- provides an
  exact, non-conjectural occurrence of $\varphi$ in a mass spectrum: the ratio m2/m1 = $\varphi$ in the
  perturbed critical Ising model, confirmed experimentally in CoNb2O6 (Coldea et al.,
  Science 327 (2010) 177). This is the only $\varphi$-in-field-theory result treated as established
  in this paper. Second, discrete scale invariance (DSI) in systems with self-similar
  hierarchical structure produces log-periodic corrections to RG flows whose preferred
  scaling ratio can, in specific models, be constrained to algebraic numbers including $\varphi

  (Sornette, Phys. Rep. 297 (1998) 239). Third, affine Toda field theories based on
  non-crystallographic Coxeter groups H2, H3, H4 produce mass pairs differing by $\varphi
  through an exact embedding construction (Fring and Korff, Nucl. Phys. B 729 (2005)
  361).

  Two candidate $\varphi$-coupling expressions are carefully distinguished throughout: the
  algebraic power variant alpha$\varphi$,alg = $\varphi$-3/2 = (5-2)/2 approx 0.1180 and the logarithmic
  variant alpha$\varphi$,log = ln($\varphi$2)/pi approx 0.30635. These are structurally distinct objects;
  they must not be conflated, and neither is derived from first principles here. A symbolic
  renormalization group (sRG) toy model and a candidate $\varphi$-structured operator
  framework, both conjectural, are introduced as formal scaffolding. Strict falsifiability
  conditions, an evidence hierarchy, mechanism-by-mechanism falsification tests, a
  derivation ladder, and five work packages for the GOLDEN BRIDGE (Trinity-Pellis Bridge) programme are
  delineated. Paper 2 is also the theoretical-mechanism branch of the broader Trinity
  $S^{3}$AI --- Flos Aureus research programme (Zenodo 10.5281/zenodo.19227877).

  Keywords: golden ratio, E8 symmetry, integrable field theory, discrete scale invariance,
  symbolic renormalization group, Toda lattice, coupling constant compressibility, Lucas
  ring, Trinity $S^{3}$AI, falsifiability, $\varphi$-structured operators, non-crystallographic Coxeter
  groups